% !TeX program = lualatex
% Compile twice: lualatex 105.tex
% All references and prose are embedded; no BibTeX database or external inputs.
\documentclass[11pt,a4paper]{article}
\usepackage[margin=24mm,headheight=14pt]{geometry}
\usepackage{fontspec}
\setmainfont{Latin Modern Roman}
\setsansfont{Latin Modern Sans}
\setmonofont{Latin Modern Mono}
\usepackage{amsmath,amssymb,mathtools}
\usepackage{unicode-math}
\setmathfont{Latin Modern Math}
\usepackage{microtype}
\usepackage{enumitem,xurl,needspace}
\usepackage[dvipsnames]{xcolor}
\usepackage{hyperref}
\hypersetup{unicode=true,colorlinks=true,linkcolor=MidnightBlue,urlcolor=MidnightBlue,
 pdftitle={Research attempt: Problem 105},pdfauthor={Research notebook},bookmarksnumbered=true}
\usepackage{bookmark}
\usepackage{tocloft}
\setlength{\cftsecnumwidth}{3em}
\cftsetpnumwidth{2.5em}
\cftsetrmarg{4em}
\usepackage{titlesec}
\titleformat{\section}{\Large\bfseries\raggedright}{\thesection}{0.7em}{}
\usepackage{fancyhdr}
\pagestyle{fancy}\fancyhf{}
\fancyhead[L]{\small\sffamily Open Problems Math | Research notebook}
\fancyhead[R]{\small Problem 105 | 27 September 2026}
\fancyfoot[C]{\thepage}
\setlength{\parindent}{0pt}
\setlength{\parskip}{3pt plus 1pt}
\setlength{\emergencystretch}{3em}
\setcounter{tocdepth}{1}
\setcounter{secnumdepth}{1}
\setlist[itemize]{leftmargin=3em,itemsep=5pt,topsep=4pt,parsep=0pt}
\allowdisplaybreaks
\newcommand{\N}{\mathbb{N}}
\newcommand{\Z}{\mathbb{Z}}
\newcommand{\Q}{\mathbb{Q}}
\newcommand{\R}{\mathbb{R}}
\newcommand{\C}{\mathbb{C}}
\newcommand{\F}{\mathbb{F}}
\newcommand{\A}{\mathbb{A}}
\newcommand{\PP}{\mathbb P}
\newcommand{\E}{\mathbb{E}}
\newcommand{\eps}{\varepsilon}
\newcommand{\dd}{\,\mathrm d}
\newcommand{\sref}[2]{\hyperlink{source:#1:#2}{[S#2]}}
\newcommand{\eref}[2]{\hyperlink{extra:#1:#2}{[E#2]}}
% Turn the next switch off for a shorter reading copy. The quotations preserve
% every exactTarget field, including qualifications omitted from a short summary.
\newif\ifshowcatalogue
\showcataloguetrue
\newcommand{\cataloguescope}[1]{\ifshowcatalogue
 \paragraph{Exact scope in the supplied catalogue.}
 \begin{quote}\small #1\end{quote}\fi}
\usepackage{etoolbox}
\AtBeginEnvironment{itemize}{\small}
\numberwithin{equation}{section}
% Standalone export. Original section 105 preserved verbatim outside marked additions.
\begin{document}
\setcounter{section}{104}
% ================================================================
% problemId: problem.top500-omission-w050054-classical-szpiro-discriminantconductor-conjecture-over-q
% source releaseRank: 105; source releaseStatus: open_with_solved_subcases
% exactTarget SHA-256: bbeaaf7b1e4bbca8186beed90ee92002724a29c26c963d20de8d8cf894bbc824
\clearpage
\section{\texorpdfstring{Classical Szpiro Discriminant–Conductor Conjecture over Q}{Classical Szpiro Discriminant–Conductor Conjecture over Q}}\label{105}
\subsection{Definitions and mathematical statement}
For an elliptic curve $E/\Q$, $\Delta_E\in\Z\setminus\{0\}$ is the discriminant of a global minimal Weierstrass equation. Its conductor $N_E=\prod_p p^{f_p}$ records the local conductor exponents of its Tate-module representation. The conjecture is
\[
 \forall\varepsilon>0\ \exists C_\varepsilon>0\ \forall E/\Q:\qquad
 |\Delta_E|\le C_\varepsilon N_E^{6+\varepsilon}.
\]
The constant is independent of the curve. The selected exponent is the classical $6+\varepsilon$ conductor--discriminant bound, not a differently normalized modified Szpiro assertion.
\par\smallskip\textit{Formulation and concept references:} \sref{105}{1}.
\cataloguescope{Prove or disprove that for every epsilon>0 there is C(epsilon)>0 such that every elliptic curve E over Q with minimal discriminant Delta\_E and conductor N\_E satisfies |Delta\_E| <= C(epsilon)*N\_E\textasciicircum{}(6+epsilon).}
\subsection{Short English statement}
Can every elliptic curve's minimal discriminant be bounded uniformly by essentially the sixth power of its conductor?
% BEGIN RESEARCH ADDITION 105
\subsection{Research attempt}

\paragraph{Current frontier.}
Here $\operatorname{num}(q)$ and $\operatorname{den}(q)$ denote the absolute
numerator and positive denominator in reduced form, following \sref{105}{1}.
Conditions involving $\log\operatorname{num}(j_E)$ below apply only where
that expression is defined; $j_E=0$ is covered separately by the integral-$j$
case.
Pasten's small-denominator theorem gives
$|\Delta_E|\leq256A N_E^6\log N_E$ when
$\operatorname{den}(j_E)\leq A\log\operatorname{num}(j_E)$, with $A$ fixed.
Its local input also yields $|\Delta_E|\leq N_E^5$ for integral $j_E$.
These are uniform results within their stated subclasses, not a bound for
all rational $j$-invariants. \sref{105}{1}

A newer manuscript, posted on 15 September 2026, states the general estimate
$\log|\Delta_E|\ll N_E\log\log N_E$ and the estimate
$\log|\Delta_E|\ll_S N_E$ for curves semistable outside a fixed finite set $S$.
These are recent preprint results, not independently audited here. Even if
accepted, they remain far from the selected bound: its logarithmic form is
$\log|\Delta_E|\leq(6+\varepsilon)\log N_E+O_\varepsilon(1)$.
In particular, a bound with an unspecified exponent on $N_E$ would not settle
this exact entry. \eref{105}{1}

\paragraph{Attack route.}
Three routes are tested. One can try to control the denominator of $j$ using
its radical, but the numerator and other bad primes may be essential. One
can seek counterexamples among quadratic twists, where minimal discriminants
and conductors are tractable. Finally, local reduction data may isolate the
part of the denominator that must be paid for by other bad primes. I use
the third route to formulate a uniform sufficient estimate, and use the
first two as adversarial tests. All constants introduced below are independent
of the curve unless a subscript explicitly says otherwise.

\paragraph{Attempt.}
\textbf{1. Keeping the three local types separate.}
Write $\delta_p=v_p(\Delta_E)$, $f_p=v_p(N_E)$ and
$B=\operatorname{den}(j_E)$. At $j_E=0$, take $B=1$ and interpret
$v_p(j_E)=+\infty$. Partition the bad primes into the following disjoint sets:
\[
 \mathcal A=\{p:v_p(j_E)\geq0\},\quad
 \mathcal M=\{p:v_p(j_E)<0,\ E\text{ multiplicative at }p\},
\]
\[
 \mathcal V=\{p:v_p(j_E)<0,\ E\text{ additive at }p\}.
\]
Define the conductor factors
\begin{equation}\label{105:factors}
 K=\prod_{p\in\mathcal A}p^{f_p},\qquad
 U=\prod_{p\in\mathcal M}p,\qquad
 V=\prod_{p\in\mathcal V}p^{f_p}.
\end{equation}
Thus $N_E=KUV$, and the factors have disjoint prime support. Good primes
cannot occur in $B$. The local estimates proved in Section 3 of
\eref{105}{2} are
\begin{align}
 p\in\mathcal A &: \quad\delta_p\leq5f_p,\nonumber\\
 p\in\mathcal M &: \quad f_p=1,\quad\delta_p=-v_p(j_E),\label{105:local}\\
 p\in\mathcal V &: \quad
 \delta_p\leq3f_p-v_p(j_E)+8\,\mathbf1_{p=2}.\nonumber
\end{align}
The last term is retained: treating the prime $2$ as a tame prime would be
an unjustified simplification.

Multiplying \eqref{105:local}, without replacing each conductor factor by
$N_E$, proves
\begin{equation}\label{105:weighted}
 |\Delta_E|\leq 256\,B K^5V^3,
 \qquad
 \frac{|\Delta_E|}{N_E^{6+\varepsilon}}
 \leq256\frac{B}{U^6V^3K\,N_E^\varepsilon}.
\end{equation}
This is the useful bookkeeping lemma of the attempt. In the tame
potentially multiplicative case $I_n^*$, one has
$\delta_p=n+6$, $f_p=2$, and $-v_p(j_E)=n$; hence
\[
 p^{\delta_p-6f_p}=p^{n-6}.
\]
For multiplicative type $I_n$ the same expression is $p^{n-6}$.
The extra discriminant power from the star is exactly cancelled by the
extra conductor power at exponent six. This explains why counting only the
number of multiplicative primes loses important structure.
\sref{105}{1}

\medskip
\textbf{2. A sufficient global arithmetic estimate.}
Consider the following statement, which is \emph{not} proved here:
\begin{quote}
For every $\varepsilon>0$ there is $A_\varepsilon$ such that, for every
elliptic curve over $\Q$,
\begin{equation}\label{105:missing}
 \operatorname{den}(j_E)\leq
 A_\varepsilon U^6V^3K\,N_E^\varepsilon,
\end{equation}
with $K,U,V$ defined by \eqref{105:factors}.
\end{quote}
Combining it with \eqref{105:weighted} would give the original conjecture,
with $C_\varepsilon=256A_\varepsilon$. This is a stronger sufficient
arithmetic statement, not an asserted equivalent formulation: the first
line of \eqref{105:local} can have substantial slack. Any proposed
counterexample to \eqref{105:missing} must still be checked against the
\emph{actual} minimal discriminant before being called a counterexample to
Szpiro.

Conversely, an unbounded sequence of the original Szpiro ratios necessarily
has unbounded right-hand ratios in \eqref{105:weighted}. The latter are
therefore a legitimate filter for a counterexample search, although not a
certificate of success. This distinction prevents an upper-bound estimate
from being used backwards.

\medskip
\textbf{3. Why a denominator-only radical bound is false.}
A tempting simplification would be
\begin{equation}\label{105:false}
 B\leq C_\varepsilon\operatorname{rad}(B)^{6+\varepsilon}
 \quad\text{for every realized rational $j$-invariant.}
\end{equation}
It fails decisively. Here is an explicit realization, rather than an appeal
to a moduli-space existence statement. For coprime integers $u,v$ with
$v>0$, $u\neq0$ and $w=u-1728v\neq0$, set
\[
 E_{u,v}:\quad y^2=x^3-3uwx-2uw^2.
\]
A direct invariant calculation gives
\begin{align}
 c_4&=144uw, &c_6&=1728uw^2,\
 \Delta&=1728^2vu^2w^3, &j&=\frac uv. \label{105:jmodel}
\end{align}
Indeed $c_4^3-c_6^2=1728\Delta$, and the displayed nonzero discriminant
ensures nonsingularity. This is an integral model, \emph{not necessarily a
minimal model}; only its $j$-invariant is used in the present argument.

Take $u=1$, $v=q^m$, with a fixed prime $q\geq5$. Then
$B=q^m$ while $\operatorname{rad}(B)=q$. Letting $m\to\infty$ disproves
\eqref{105:false} for every proposed fixed constant. It does \emph{not}
disprove Szpiro: primes dividing $1-1728q^m$ also enter the displayed model,
and their contributions to the minimal conductor cannot simply be erased.
This family pinpoints the missing interaction between the denominator and
other bad primes in \eqref{105:missing}.

\medskip
\textbf{4. Counterexample search through quadratic twisting.}
Fix an elliptic curve $E/\Q$. Let $d$ be a positive squarefree integer with
\[
 d\equiv1\pmod{24N_E}.
\]
In particular $d$ is coprime to $6N_E$ and is a local square at $2$, $3$ and
all bad primes of $E$: use the square criterion modulo $8$ at $2$ and Hensel
lifting at odd primes. Let $E^{(d)}$ be the quadratic twist, with short-model
coefficients multiplied by $d^2,d^3$. Then
\begin{equation}\label{105:twist}
 |\Delta_{E^{(d)}}|=d^6|\Delta_E|,
 \qquad N_{E^{(d)}}=d^2N_E.
\end{equation}

\emph{Local verification.} At the old bad primes the curves are locally
isomorphic, so their local minimal discriminants and conductors agree.
The same holds at $2$ and $3$. At a prime $p\mid d$, necessarily $p\geq5$,
start with a minimal short equation for $E$ with unit discriminant. The twist
has discriminant valuation $6$. It is minimal because lowering that valuation
would require subtracting a positive multiple of $12$ from it. The original
Tate-module representation is unramified; after the ramified quadratic twist,
tame inertia acts by $-I$. Its invariant space has dimension zero and its
Swan conductor is zero, so its conductor exponent is $2$. All remaining
primes stay good. This proves both equalities in \eqref{105:twist}.
The discriminant-scaling and conductor conventions are those in the
source's local discussion. \sref{105}{1}

It follows exactly, not just asymptotically, that
\begin{equation}\label{105:twist-ratio}
 \frac{|\Delta_{E^{(d)}}|}{N_{E^{(d)}}^{6+\varepsilon}}
 =\frac{|\Delta_E|}{N_E^{6+\varepsilon}}\,d^{-6-2\varepsilon}.
\end{equation}
Thus this natural unbounded twist search moves \emph{away} from a
counterexample. In \eqref{105:factors}, these new primes instead multiply
$K$ by $d^2$ while leaving $B,U,V$ unchanged. The sufficient-estimate ratio
in \eqref{105:weighted} falls only as $d^{-2-2\varepsilon}$, displaying
explicitly the slack between that criterion and the actual Szpiro ratio.
A constant depending on the initially fixed curve cannot be promoted to a
constant uniform over all curves.

\medskip
\textbf{5. Combining the small-denominator route with the recent height input.}
There is a modest, conditional sharpening that can be made without pretending
to solve the denominator problem. The local divisibility bound
$|\Delta_E|\leq16 B N_E^5$ and the comparison
$h(j_E)\ll1+h(E)$ are given in \sref{105}{1}.
If one accepts the new preprint's estimate
$h(E)\ll N_E\log\log(3N_E)$, then for fixed $A,\beta>0$,
\[
 B\leq A(\log\operatorname{num}(j_E))^\beta
 \quad\Longrightarrow\quad
 |\Delta_E|\ll_\beta A\,
 N_E^{5+\beta}\bigl(\log\log(3N_E)\bigr)^\beta.
\]
This follows simply by bounding $\log\operatorname{num}(j_E)$ by $h(j_E)$
and substituting. At $\beta=1$ the logarithmic loss in the older subclass
bound becomes an iterated logarithm. The dependence on the recent input is
explicit; no independent verification of that input, and no novelty claim
for this combination, is made. \eref{105}{1}

\paragraph{Stress test.}
The argument has kept minimal discriminants distinct from discriminants of
convenient equations; \eqref{105:jmodel} is used only to realize $j$.
The wild-prime allowance in \eqref{105:local} is not absorbed into an
unmentioned curve-dependent constant. The twisting calculation imposes a
congruence that makes the local behavior at the old bad primes explicit.
All powers of $N_E$ in \eqref{105:weighted} have been accounted for before
formulating the missing estimate.

Neither a finite set of large Szpiro ratios nor the false estimate
\eqref{105:false} provides a counterexample to the quantified conjecture:
for a fixed $\varepsilon$ one needs ratios unbounded over curves. The symbolic
checks in the companion script verify \eqref{105:jmodel} and the exponent
cancellations. They do not compute minimal models or test
\eqref{105:missing} over all elliptic curves. In particular, no invocation
of an unproved $abc$ or modified Szpiro statement has been hidden in the
local-to-global multiplication.

\paragraph{Outcome.}
\textbf{Outcome: Conditional reduction.}

The proved result of this attempt is the weighted inequality
\eqref{105:weighted}, obtained by retaining the distinct local conductor
factors, together with the explicit failure of the denominator-only radical
lemma and the exact negative result for the specified twist search.
Estimate \eqref{105:missing}, uniformly for every curve and every
$\varepsilon>0$, would imply the exact classical Szpiro statement. It is
not established, and is stronger than the original conjecture at some local
types. The recent-height combination concerns only a small-denominator
subclass. No complete proof or counterexample candidate is claimed.
% END RESEARCH ADDITION 105
\subsection{Sources}
\begin{itemize}\raggedright
\item[\textnormal{[S1]}]\hypertarget{source:105:1}{} H. Pasten, Szpiro's conjecture when the denominator of the j-invariant is small, arXiv v2 (15 August 2023).
\url{https://arxiv.org/html/2308.07114}.
{\small\textit{Catalogue formulation source.}}
% BEGIN REFERENCE ADDITION 105
\item[\textnormal{[E1]}]\hypertarget{extra:105:1}{} H. Pasten,
\emph{Improved bounds for Szpiro's conjecture}, arXiv:2609.17390v1
(15 September 2026), Theorem 1.1, Corollaries 1.2--1.3,
and equation (1.1).
\url{https://arxiv.org/html/2609.17390v1}.
{\small\textit{Consulted 27 September 2026; recent preprint, not independently proof-audited here.}}
\item[\textnormal{[E2]}]\hypertarget{extra:105:2}{} H. Pasten,
\emph{Szpiro's conjecture when the denominator of the $j$-invariant is small},
arXiv:2308.07114v2 (15 August 2023), Theorem 1.2, Corollary 1.3,
Section 2, Lemmas 3.1--3.3 and Corollary 3.4.
\url{https://arxiv.org/html/2308.07114v2}.
{\small\textit{Same work as [S1]; this added locator records the exact results checked on 27 September 2026. The HTML's later generated date is not treated as a new arXiv revision.}}
% END REFERENCE ADDITION 105
\end{itemize}

\end{document}
